复共轭与换叶分别作用为
\begin{equation}\label{eq:symmetry-domains}
z\mapsto\conj{z}:D_\pm\longrightarrow D_\mp,\qquad
h:z\mapsto-\dfrac{q_0^2}{z}:D_\pm\longrightarrow D_\mp.
\end{equation}
二者对易, 其复合保持 $D_\pm$. 以下先在普通连续谱上比较 Jost 归一化, 再将相应列关系延拓到解析域.

\begin{proposition}[Jost 解的共轭与换叶约化]\label{prop:symmetry}
令
\[
\sigma_2=\begin{pmatrix}0&-\ii\\\ii&0\end{pmatrix},\qquad
C(z)=\dfrac{\ii z}{q_0}\sigma_1.
\]
在定理 \ref{thm:jost} 的条件下,
\begin{equation}\label{eq:conjugation}
\phi_\pm(z)=\sigma_2\conj{\phi_\pm(\conj{z})}\sigma_2,\qquad
\mu_\pm(z)=\sigma_2\conj{\mu_\pm(\conj{z})}\sigma_2,
\end{equation}
\begin{equation}\label{eq:musheet}
\phi_\pm(h(z))=\phi_\pm(z)C(z),\qquad
\mu_\pm(h(z))=\mu_\pm(z)C(z).
\end{equation}
逐列写为
\begin{align}
\phi_{\pm,1}(z)&=\ii\sigma_2\conj{\phi_{\pm,2}(\conj{z})},&
\phi_{\pm,2}(z)&=-\ii\sigma_2\conj{\phi_{\pm,1}(\conj{z})},\label{eq:jost-conjugate-columns}\\
\phi_{\pm,1}(h(z))&=\dfrac{\ii z}{q_0}\phi_{\pm,2}(z),&
\phi_{\pm,2}(h(z))&=\dfrac{\ii z}{q_0}\phi_{\pm,1}(z).\label{eq:jost-sheet-columns}
\end{align}
修正列 $\mu_{\pm,j}$ 满足相同关系.
\end{proposition}
\begin{proof}
由 $\sigma_2\sg\sigma_2=-\sg$、$\sigma_2\conj{Q}\sigma_2=Q$ 及 $k(\conj{z})=\conj{k(z)}$,
\[
U(z)=\sigma_2\conj{U(\conj{z})}\sigma_2,\qquad
Y(z)=\sigma_2\conj{Y(\conj{z})}\sigma_2.
\]
又有 $\theta(\conj{z})=\conj{\theta(z)}$, 故
\[
\sigma_2\conj{Y(\conj{z})e^{\ii\theta(\conj{z})\sg}}\sigma_2
=Y(z)e^{\ii\theta(z)\sg}.
\]
共轭后的 Jost 解满足同一空间方程和同一归一化, 定理 \ref{thm:jost} 的唯一性给出 \eqref{eq:conjugation}.

为逐元检查, 有
\[
\sigma_2\sg\sigma_2
=\begin{pmatrix}-1&0\\0&1\end{pmatrix}=-\sg,
\qquad
\sigma_2\sigma_1\sigma_2=-\sigma_1.
\]
且
\[
\begin{aligned}
\sigma_2\conj{Y(\conj{z})}\sigma_2
&=\sigma_2\left(I+\dfrac{\ii q_0}{z}\sigma_1\right)\sigma_2
=I-\dfrac{\ii q_0}{z}\sigma_1=Y(z),\\
\sigma_2\conj{U(\conj{z})}\sigma_2
&=\sigma_2(\ii k\sg+\conj Q)\sigma_2
=-\ii k\sg+Q=U(z).
\end{aligned}
\]
因此对 $\phi_\pm$ 复共轭并左右乘 $\sigma_2$,
得到的函数既满足同一 $x$ 方程, 又有同一端点归一化.
对矩阵取右端的两列便得到
$\phi_{\pm,1}(z)=\ii\sigma_2\conj{\phi_{\pm,2}(\conj{z})}$、
$\phi_{\pm,2}(z)=-\ii\sigma_2\conj{\phi_{\pm,1}(\conj{z})}$.

 在传播方程适用时, $\sigma_2\conj{R}\sigma_2=-R$ 同时保证 $V(z)=\sigma_2\conj{V(\conj{z})}\sigma_2$.

换叶保持 $k$ 而使 $\lambda,\theta$ 变号. 直接由 \eqref{eq:E} 计算,
\begin{equation}\label{eq:Esheet}
Y(h(z))=I+\dfrac{\ii z}{q_0}\sigma_1=Y(z)C(z),\qquad
C(z)\sg=-\sg C(z).
\end{equation}
于是 $\phi_\pm(z)C(z)$ 在 $h(z)$ 处满足同一空间方程, 且
\[
Y(z)e^{\ii\theta(z)\sg}C(z)
=Y(h(z))e^{\ii\theta(h(z))\sg}.
\]
再次由唯一性得到 $\phi_\pm$ 的换叶关系. 去除指数时使用
$C(z)e^{\ii\theta(z)\sg}=e^{-\ii\theta(z)\sg}C(z)$,
便得到 $\mu_\pm$ 的关系.

换叶关系中的常数矩阵也可直接相乘:
\[
\begin{aligned}
Y(z)C(z)
&=\left(I-\dfrac{\ii q_0}{z}\sigma_1\right)
\dfrac{\ii z}{q_0}\sigma_1\\
&=\dfrac{\ii z}{q_0}\sigma_1+I
=I-\dfrac{\ii q_0}{h(z)}\sigma_1=Y(h(z)).
\end{aligned}
\]
这里用了 $h(z)=-q_0^2/z$, 因而
$-\ii q_0/h(z)=\ii z/q_0$.
再由 $\sigma_1\sg=-\sg\sigma_1$,
有
$e^{\ii\theta\sg}C=C e^{-\ii\theta\sg}$,
从而两种 Jost 归一化在换叶下严格一致.

 右乘矩阵的两列分别给出 \eqref{eq:jost-conjugate-columns}--\eqref{eq:jost-sheet-columns}. 变换后的列具有 \eqref{eq:symmetry-domains} 所确定的解析域, 因此同一 Volterra 唯一性论证在各解析域内成立.
\end{proof}

\begin{proposition}[散射矩阵与反射数据的约化]\label{prop:scattering-symmetry}
散射矩阵满足
\begin{equation}\label{eq:scattering-symmetry}
S(z)=\sigma_2\conj{S(\conj{z})}\sigma_2,\qquad
S(h(z))=\sigma_1S(z)\sigma_1.
\end{equation}
特别地,
\begin{equation}\label{eq:sconj}
s_{22}(z)=\conj{s_{11}(\conj{z})},\qquad
s_{12}(z)=-\conj{s_{21}(\conj{z})},
\end{equation}
\begin{equation}\label{eq:ssheet}
s_{22}(h(z))=s_{11}(z),\qquad s_{12}(h(z))=s_{21}(z),
\end{equation}
\begin{equation}\label{eq:reflection-symmetry}
\rho(z)=-\conj{\widetilde\rho(\conj{z})},\qquad
\rho(h(z))=\widetilde\rho(z).
\end{equation}
透射系数的关系在相应解析域内成立. 在实谱的普通点,
\begin{equation}\label{eq:scattering-unitarity}
S(z)^\dagger S(z)=I,\qquad
|s_{11}(z)|^2+|s_{21}(z)|^2=1.
\end{equation}
\end{proposition}
\begin{proof}
将 \eqref{eq:conjugation} 代入 $\phi_+=\phi_-S$, 并利用 $\phi_-$ 的可逆性, 得到 \eqref{eq:scattering-symmetry} 的第一式. 比较
\[
\phi_+(z)C(z)=\phi_-(z)S(z)C(z)
=\phi_-(z)C(z)S(h(z)),
\]
得到
\[
S(h(z))=C(z)^{-1}S(z)C(z)=\sigma_1S(z)\sigma_1.
\]
逐元展开即得 \eqref{eq:sconj}--\eqref{eq:ssheet}; 除以相应透射系数得到 \eqref{eq:reflection-symmetry}.

所说的逐元展开如下:
\[
\begin{aligned}
\sigma_2\conj{S(\conj{z})}\sigma_2
&=\begin{pmatrix}
\conj{s_{22}(\conj{z})}&-\conj{s_{21}(\conj{z})}\\
-\conj{s_{12}(\conj{z})}&\conj{s_{11}(\conj{z})}
\end{pmatrix},\\
\sigma_1S(z)\sigma_1
&=\begin{pmatrix}s_{22}(z)&s_{21}(z)\\
s_{12}(z)&s_{11}(z)\end{pmatrix}.
\end{aligned}
\]
例如第一种约化给出
$s_{21}(z)=-\conj{s_{12}(\conj{z})}$,
第二种给出
$s_{21}(h(z))=s_{12}(z)$.
所以
\[
\begin{aligned}
\rho(z)
&=\dfrac{s_{21}(z)}{s_{11}(z)}
=-\dfrac{\conj{s_{12}(\conj{z})}}
{\conj{s_{22}(\conj{z})}}
=-\conj{\widetilde\rho(\conj{z})},\\
\rho(h(z))
&=\dfrac{s_{21}(h(z))}{s_{11}(h(z))}
=\dfrac{s_{12}(z)}{s_{22}(z)}
=\widetilde\rho(z).
\end{aligned}
\]
这些都是相应透射系数非零处的恒等式.

 解析透射系数的关系也可直接对 \eqref{eq:wronskian} 应用 Jost 列的约化得到.

在实 $z$ 上, $k,\lambda,F$ 为实数, 且 $U^\dagger=-U$. 因而
\[
\partial_x(\phi_\pm^\dagger\phi_\pm)
=\phi_\pm^\dagger(U^\dagger+U)\phi_\pm=0,\qquad
Y^\dagger Y=\gamma I.
\]
振荡归一化给出 $\phi_\pm^\dagger\phi_\pm=\gamma I$. 代入散射关系并除以 $\gamma>0$, 即得 \eqref{eq:scattering-unitarity}.

在实轴上具体有
\[
Y^\dagger
=\begin{pmatrix}1&\dfrac{\ii q_0}{z}\\
\dfrac{\ii q_0}{z}&1\end{pmatrix},
\qquad
Y^\dagger Y
=\begin{pmatrix}1+\dfrac{q_0^2}{z^2}&0\\
0&1+\dfrac{q_0^2}{z^2}\end{pmatrix}
=\gamma I.
\]
因此
\[
\gamma I=\phi_+^\dagger\phi_+
=S^\dagger\phi_-^\dagger\phi_-S
=\gamma S^\dagger S,
\]
即 $S^\dagger S=I$. 取其 $(1,1)$ 元得到
$\conj{s_{11}}s_{11}+\conj{s_{21}}s_{21}=1$.


\end{proof}
